% Présentation de soutenance (Beamer, 16:9) -- projet de démonstration
% Thème personnalisé : palette crème / bordeaux, barre de progression, figures TikZ et pgfplots.
% Compilation : latexmk -pdf main.tex
\documentclass[aspectratio=169,11pt]{beamer}

\usepackage[T1]{fontenc}
\usepackage[utf8]{inputenc}
\usepackage[french]{babel}
\usepackage{amsmath,amssymb}
\usepackage{newpxtext,newpxmath}
\usepackage{microtype}
\usepackage{booktabs}
\usepackage{tikz}
\usetikzlibrary{arrows.meta,calc,decorations.pathreplacing}
\usepackage{pgfplots}
\pgfplotsset{compat=1.18}

% ------------------------------------------------------------------ palette
\definecolor{wine}{HTML}{7F1D1D}
\definecolor{ink}{HTML}{1C1917}
\definecolor{sea}{HTML}{1D4E89}
\definecolor{moss}{HTML}{2F6B3A}
\definecolor{ochre}{HTML}{B7791F}
\definecolor{paper}{HTML}{FBF7EE}
\definecolor{sand}{HTML}{F3ECDA}
\definecolor{line}{HTML}{DDD2B8}

% ------------------------------------------------------------------ thème
\usefonttheme{serif}
\setbeamertemplate{navigation symbols}{}
\setbeamercolor{normal text}{fg=ink,bg=paper}
\setbeamercolor{structure}{fg=wine}
\setbeamercolor{frametitle}{fg=wine}
\setbeamercolor{title}{fg=ink}
\setbeamercolor{item}{fg=wine}
\setbeamercolor{block title}{fg=wine,bg=sand}
\setbeamercolor{block body}{fg=ink,bg=sand!55}
\setbeamerfont{frametitle}{series=\bfseries,size=\Large}
\setbeamerfont{block title}{series=\bfseries,size=\normalsize}
\setbeamertemplate{itemize item}{\raisebox{0.15ex}{\scriptsize$\blacksquare$}}
\setbeamertemplate{itemize subitem}{\raisebox{0.1ex}{\tiny$\blacksquare$}}
\setbeamertemplate{enumerate item}{\insertenumlabel.}
\setbeamertemplate{blocks}[rounded][shadow=false]
\setlength{\leftmargini}{1.2em}

\setbeamertemplate{frametitle}{%
  \vspace{0.9em}\nointerlineskip
  {\usebeamerfont{frametitle}\usebeamercolor[fg]{frametitle}\insertframetitle\par}
  \vspace{-0.1em}{\color{line}\rule{\textwidth}{0.8pt}}\par
}
\setbeamertemplate{headline}{}
\setbeamertemplate{footline}{%
  \hbox{\hspace{1.4em}\color{ink!60}\scriptsize Prénom Nom \textperiodcentered{} Soutenance de master \hfill
        \insertframenumber\,/\,\inserttotalframenumber\hspace{1.4em}}%
  \vskip 0.55em
  \hbox{\begin{tikzpicture}[overlay]
    \fill[line] (0,0) rectangle (\paperwidth,0.07);
    \fill[wine] (0,0) rectangle ({\paperwidth*\insertframenumber/\inserttotalframenumber},0.07);
  \end{tikzpicture}}%
}\setbeamersize{text margin left=1.4em,text margin right=1.4em}

\setbeamertemplate{title page}{%
  \begin{tikzpicture}[remember picture,overlay]
    \fill[wine] (current page.north west) rectangle ([xshift=0.45cm]current page.south west);
  \end{tikzpicture}%
  \vspace{2.2em}
  \hspace{1.2em}\begin{minipage}{0.86\textwidth}\raggedright
    {\color{wine}\scshape\small Université Exemple \textperiodcentered{} Département de Mathématiques}\par
    \vspace{1.6em}
    {\fontsize{24}{28}\selectfont\bfseries\inserttitle\par}
    \vspace{0.6em}
    {\Large\itshape\insertsubtitle\par}
    \vspace{2.2em}
    {\large\insertauthor\par}
    \vspace{0.3em}
    {\small\color{ink!70}Sous la direction de Pr.~Nom du directeur\par}
    \vspace{0.2em}
    {\small\color{ink!70}\insertdate\par}
  \end{minipage}
}

\newcommand{\R}{\mathbb{R}}
\newcommand{\norm}[1]{\lVert #1\rVert}

\pgfplotsset{
  sci labels/.style={nodes near coords style={font=\scriptsize,/pgf/number format/sci generic={mantissa sep=\times,exponent={10^{####1}}}}},
  slideplot/.style={
    font=\footnotesize, tick label style={font=\scriptsize}, label style={font=\footnotesize},
    axis line style={ink!70}, tick style={ink!70},
    grid=major, grid style={line!70, line width=0.3pt},
    legend style={font=\scriptsize, fill=paper, fill opacity=0.9, draw=line, text opacity=1, cells={anchor=west}},
    every axis plot/.append style={thick},
  },
}

\title{Méthodes numériques pour les équations différentielles ordinaires}
\subtitle{Analyse de convergence et applications}
\author{Prénom Nom}
\date{Octobre 2026}

\begin{document}

% ------------------------------------------------------------------ 1
\begin{frame}[plain,t]
  \titlepage
\end{frame}

% ------------------------------------------------------------------ 2
\begin{frame}{Plan de l'exposé}
  \begin{columns}[T]
    \column{0.52\textwidth}
    \begin{enumerate}
      \setlength{\itemsep}{0.7em}
      \item Le problème de Cauchy
      \item Trois méthodes à un pas
      \item Consistance, ordre et convergence
      \item Expériences numériques
      \item Application : proies et prédateurs
      \item Conclusion et perspectives
    \end{enumerate}
    \column{0.42\textwidth}
    \begin{block}{Question centrale}
      Comment relier la précision d'un pas de calcul à l'erreur accumulée sur tout l'intervalle ?
    \end{block}
  \end{columns}
\end{frame}

% ------------------------------------------------------------------ 3
\begin{frame}{Le problème de Cauchy}
  \begin{block}{Problème}
    \[
      y'(t)=f\bigl(t,y(t)\bigr),\qquad y(t_0)=y_0,\qquad t\in[t_0,T]
    \]
  \end{block}
  \vspace{0.4em}
  \begin{columns}[T]
    \column{0.52\textwidth}
    \begin{block}{Hypothèse}
      $f$ continue et lipschitzienne en $y$ :
      \[
        \norm{f(t,u)-f(t,v)}\le L\,\norm{u-v}
      \]
    \end{block}
    \column{0.44\textwidth}
    \begin{block}{Théorème (Cauchy--Lipschitz)}
      Le problème admet une \emph{unique} solution sur $[t_0,T]$.
    \end{block}
  \end{columns}
  \vspace{0.5em}
  \textbf{Problème test :} $y'=y$, $y(0)=1$, de solution exacte $y(t)=e^{t}$.
\end{frame}

% ------------------------------------------------------------------ 4
\begin{frame}{Trois méthodes à un pas}
  Schéma général : \quad $y_{n+1}=y_n+h\,\Phi(t_n,y_n;h)$, \quad $t_n=t_0+nh$.
  \vspace{0.8em}
  \begin{columns}[T]
    \column{0.31\textwidth}
    \begin{block}{Euler}
      \small $\Phi=f(t,y)$\\[0.4em]
      1 évaluation par pas\\
      ordre $1$
    \end{block}
    \column{0.31\textwidth}
    \begin{block}{Heun}
      \small $\Phi=\tfrac12\bigl[f(t,y)+f(t+h,\,y+hf(t,y))\bigr]$\\[0.4em]
      2 évaluations par pas\\
      ordre $2$
    \end{block}
    \column{0.31\textwidth}
    \begin{block}{Runge--Kutta 4}
      \small $\Phi=\tfrac16\,(k_1+2k_2+2k_3+k_4)$\\[0.4em]
      4 évaluations par pas\\
      ordre $4$
    \end{block}
  \end{columns}
\end{frame}

% ------------------------------------------------------------------ 5
\begin{frame}{Erreur locale et erreur globale}
  \begin{columns}[c]
    \column{0.52\textwidth}
    \centering
    \begin{tikzpicture}[x=2.1cm,y=0.62cm,>={Stealth[length=2mm]}]
      \draw[->,ink!60] (0,0) -- (3.5,0) node[right] {\scriptsize $t$};
      \draw[->,ink!60] (0,0) -- (0,7.1) node[above] {\scriptsize $y$};
      \draw[thick,sea,domain=0:3.3,samples=50,smooth,variable=\x] plot ({\x},{1+0.6*\x+0.35*\x*\x});
      \draw[ochre,dashed,thick] (0.35,{1.95+1.3*(0.35-1)}) -- (1,1.95);
      \draw[wine,very thick] (1,1.95) -- (2.2,3.51);
      \draw[ink!50,dotted] (1,0) -- (1,1.95);
      \draw[ink!50,dotted] (2.2,0) -- (2.2,3.51);
      \node[below,font=\scriptsize] at (1,0) {$t_n$};
      \node[below,font=\scriptsize] at (2.2,0) {$t_{n+1}$};
      \fill[sea] (1,1.95) circle (2pt) node[above left,font=\scriptsize] {$y(t_n)=y_n$};
      \fill[wine] (2.2,3.51) circle (2pt) node[below right,font=\scriptsize] {$y_{n+1}$};
      \fill[sea] (2.2,4.014) circle (2pt) node[above left,font=\scriptsize] {$y(t_{n+1})$};
      \draw[decorate,decoration={brace,amplitude=3pt,mirror},wine] (2.2,3.51) -- (2.2,4.014);
      \node[wine,font=\scriptsize,anchor=west] at (2.4,3.76) {erreur locale};
    \end{tikzpicture}
    \column{0.43\textwidth}
    \begin{block}{Erreur locale}
      $\delta(t;h)=y(t+h)-y(t)-h\,\Phi\bigl(t,y(t);h\bigr)$
    \end{block}
    \begin{block}{Erreur globale}
      $e_n=y(t_n)-y_n$
    \end{block}
    \small Consistance d'ordre $p$ : $\norm{\delta}\le C\,h^{p+1}$.
  \end{columns}
\end{frame}

% ------------------------------------------------------------------ 6
\begin{frame}{Théorème de convergence}
  \begin{block}{Théorème}
    Si la méthode est consistante d'ordre $p$ et si $\Phi$ est $\Lambda$-lipschitzienne en $y$, alors
    \[
      \norm{e_n}\;\le\;\frac{C}{\Lambda}\Bigl(e^{\Lambda(t_n-t_0)}-1\Bigr)\,h^{p}.
    \]
  \end{block}
  \vspace{0.3em}
  \textbf{Idée de la preuve}
  \begin{enumerate}
    \setlength{\itemsep}{0.35em}
    \item $e_{n+1}=e_n+h\bigl[\Phi(t_n,y(t_n);h)-\Phi(t_n,y_n;h)\bigr]+\delta_n$
    \item $\norm{e_{n+1}}\le(1+h\Lambda)\norm{e_n}+C\,h^{p+1}$
    \item somme géométrique et $(1+h\Lambda)^n\le e^{\Lambda(t_n-t_0)}$
  \end{enumerate}
  \vspace{0.2em}
  \small Conséquence : une méthode consistante d'ordre $p$ converge à l'ordre $p$.
\end{frame}

% ------------------------------------------------------------------ 7
\begin{frame}{Les ordres observés confirment la théorie}
  \begin{columns}[c]
    \column{0.58\textwidth}
    \begin{tikzpicture}
    \begin{loglogaxis}[
      slideplot, width=\linewidth, height=5.4cm,
      xlabel={pas $h$}, ylabel={erreur en $t=1$},
      xmin=3e-3, xmax=0.4, ymin=1e-14, ymax=1, legend pos=south east,
      xtick={0.004,0.0156,0.0625,0.25}, xticklabels={$2^{-8}$,$2^{-6}$,$2^{-4}$,$2^{-2}$},
      ytick={1e-12,1e-9,1e-6,1e-3,1},
    ]
      \addplot[wine,mark=square*,mark size=1.6pt] table[col sep=comma,x=h,y=euler] {convergence.csv};
      \addlegendentry{Euler}
      \addplot[sea,mark=*,mark size=1.6pt] table[col sep=comma,x=h,y=heun] {convergence.csv};
      \addlegendentry{Heun}
      \addplot[moss,mark=triangle*,mark size=2pt] table[col sep=comma,x=h,y=rk4] {convergence.csv};
      \addlegendentry{RK4}
    \end{loglogaxis}
    \end{tikzpicture}
    \column{0.38\textwidth}
    \small
    Sur $y'=y$, $y(0)=1$ :
    \begin{itemize}
      \setlength{\itemsep}{0.4em}
      \item pente $1$ pour Euler
      \item pente $2$ pour Heun
      \item pente $4$ pour RK4
    \end{itemize}
    \vspace{0.3em}
    Diviser $h$ par $2$ divise l'erreur par $2$, $4$ ou $16$.
  \end{columns}
\end{frame}

% ------------------------------------------------------------------ 8
\begin{frame}{Le coût de la précision}
  \begin{columns}[c]
    \column{0.55\textwidth}
    \begin{tikzpicture}
    \begin{axis}[
      slideplot, width=\linewidth, height=5.4cm,
      ybar, bar width=24pt, ymode=log, log origin=infty,
      symbolic x coords={Euler,Heun,RK4}, xtick=data,
      ylabel={évaluations de $f$}, ymin=10, ymax=1e8, enlarge x limits=0.3,
      point meta=explicit,
      nodes near coords={\pgfmathprintnumber[sci,precision=1]{\pgfplotspointmeta}},
      sci labels,
      ytick={10,1e2,1e4,1e6,1e8},
    ]
      \addplot[fill=ochre!70,draw=ochre!90!ink] table[col sep=comma,x=method,y=evals,meta=evals] {cost.csv};
    \end{axis}
    \end{tikzpicture}
    \column{0.40\textwidth}
    \begin{block}{Erreur $\le 10^{-6}$ en $t=1$}
      \small
      \begin{tabular}{@{}lrr@{}}
        \toprule
        & pas & évaluations\\
        \midrule
        Euler & 1\,359\,469 & 1\,359\,469\\
        Heun  & 673         & 1\,346\\
        RK4   & 13          & 52\\
        \bottomrule
      \end{tabular}
    \end{block}
    \small RK4 coûte plus de $10^4$ fois moins qu'Euler.
  \end{columns}
\end{frame}

% ------------------------------------------------------------------ 9
\begin{frame}{Application : proies et prédateurs}
  \begin{columns}[T]
    \column{0.30\textwidth}
    \small
    Modèle de Lotka--Volterra :
    \[
      \begin{cases}
        x'=\alpha x-\beta xy\\
        y'=\delta xy-\gamma y
      \end{cases}
    \]
    $\alpha=1$, $\beta=0{,}1$, $\gamma=1{,}5$, $\delta=0{,}075$\\
    $(x_0,y_0)=(10,5)$
    \column{0.34\textwidth}
    \begin{tikzpicture}
    \begin{axis}[
      slideplot, width=\linewidth, height=5.2cm,
      xlabel={temps $t$}, ylabel={population}, xmin=0, xmax=30, ymin=0, legend pos=north east,
    ]
      \addplot[sea,no marks] table[col sep=comma,x=t,y=prey] {predator_prey.csv};
      \addlegendentry{proies}
      \addplot[wine,no marks] table[col sep=comma,x=t,y=predator] {predator_prey.csv};
      \addlegendentry{prédateurs}
    \end{axis}
    \end{tikzpicture}
    \column{0.30\textwidth}
    \begin{tikzpicture}
    \begin{axis}[
      slideplot, width=\linewidth, height=5.2cm,
      xlabel={proies}, ylabel={prédateurs}, xmin=0, ymin=0, legend pos=north east,
    ]
      \addplot[ochre,no marks] table[col sep=comma,x=prey,y=predator] {predator_prey.csv};
      \addlegendentry{orbite}
      \addplot[only marks,mark=*,mark size=2.2pt,wine] coordinates {(20,10)};
      \addlegendentry{équilibre}
    \end{axis}
    \end{tikzpicture}
  \end{columns}
  \vspace{0.2em}
  \small Oscillations de période $\approx 5{,}5$. Intégrale première $I=\delta x-\gamma\ln x+\beta y-\alpha\ln y$ : avec RK4, elle varie de moins de $1{,}2\times10^{-5}$ sur $[0,30]$.
\end{frame}

% ------------------------------------------------------------------ 10
\begin{frame}{Conclusion et perspectives}
  \begin{columns}[T]
    \column{0.47\textwidth}
    \begin{block}{Ce qu'il faut retenir}
      \small
      \begin{itemize}
        \setlength{\itemsep}{0.35em}
        \item Consistance d'ordre $p$ $\Rightarrow$ convergence d'ordre $p$
        \item Ordres observés : $1$, $2$ et $4$
        \item Coût : RK4 très en avance à tolérance fine
        \item Invariant du modèle bien conservé
      \end{itemize}
    \end{block}
    \column{0.47\textwidth}
    \begin{block}{Perspectives}
      \small
      \begin{itemize}
        \setlength{\itemsep}{0.35em}
        \item Pas adaptatif et contrôle de l'erreur
        \item Méthodes implicites pour les problèmes raides
        \item Méthodes symplectiques
      \end{itemize}
    \end{block}
  \end{columns}
\end{frame}

% ------------------------------------------------------------------ 11
\begin{frame}[plain,c]
  \begin{tikzpicture}[remember picture,overlay]
    \fill[wine] (current page.north west) rectangle ([xshift=0.45cm]current page.south west);
  \end{tikzpicture}
  \centering
  {\fontsize{34}{40}\selectfont\bfseries Merci de votre attention}\\[1.2em]
  {\Large\itshape Questions ?}
\end{frame}

\end{document}
